% !TeX root = ../../Book5.tex
% 纲要标题：### 8.4 模表示理论补充
% 纲要位置：工作记录/计划纲要.md，第 4001 行。

\section{模表示理论补充}
\label{b5:sec:0804}


% 纲要范围：
%
% * 本章不同于第5章 Brauer 诱导理论
% * 缺陷群和块论的学习地图
% * 不把特征零分类原样搬到正特征
% * 选定诱导格的实际约化同构与 Grothendieck 群中的相容公式；相对迹与 S3 块的 Brauer 像
% * 根层相同但不同构的二维模族；循环 p 群上稳定同态、第一合冲模与短正合扩张的直接计算
%
% ---
%

本章的几个计算已经说明，特征零表示、模表示和投射模之间有联系，却不能彼此替换。本节用约化与诱导的相容性说明第一种联系，再用一个能够直接计算的截断映射说明块的缺陷群。更深的块分类还需要其他工具，这里先把问题和所需理论说明清楚。

\subsection{模表示与 Brauer 诱导理论的区别}
\label{b5:subsec:0804:01}

\chapref{b5:ch:05}的 Brauer 诱导定理在特征零表示环中，把普通特征标写成由某些子群的一维特征标诱导而来的整数线性组合。本章的 Brauer 特征标则在特征 \(p\) 中先取表示，再把 \(p\)-正则特征值提升到特征零；分解矩阵记录整格约化的合成因子。两种理论处理不同的构造，不能因为都冠以 Brauer 之名就将其结论互换。
\begin{proposition}\label{b5:prop:reduction-induction}
	设 \(G\) 为有限群，\(H\le G\)，\(p\) 为素数，并给定对 \(G,H\) 都分裂的 \(p\)-模系统 \((K,\mathcal O,k)\)。记 \(d_G:R_K(G)\to G_0(kG)\)、\(d_H:R_K(H)\to G_0(kH)\) 为分解映射，其中表示群及 Grothendieck 群\footnote{格罗滕迪克（Alexandre Grothendieck，1928-2014），法籍数学家，重塑代数几何的概形与上同调方法。}都取有限维左模。它们与限制、诱导相容：
	\[
		d_G\operatorname{Ind}_H^G =\operatorname{Ind}_H^Gd_H,\qquad d_H\operatorname{Res}_H^G =\operatorname{Res}_H^Gd_G.
	\]
	右侧的诱导、限制作用于由短正合列定义的 \(G_0\) 群。
\end{proposition}
\begin{proof}
	限制一个稳定满格仍是稳定满格，其约化显然与限制交换。若 \(L\) 是 \(KH\)-模 \(V\) 中的满格，则 \(\mathcal OG\otimes_{\mathcal OH}L\) 是 \(\operatorname{Ind}_H^G V\) 中的稳定满格：按左陪集代表，\(\mathcal OG\) 是有限自由右 \(\mathcal OH\) 模，这个张量是若干份 \(L\) 的直和，且扩到 \(K\) 后得到所需诱导模。模 \(\pi\) 后自然同构于 \(kG\otimes_{kH}(L/\pi L)\)。 诱导因右侧自由性保持正合列，限制也保持正合列，所以二者确实诱导 \(G_0\) 上的映射。结合格选择无关性得到公式。
\end{proof}
证明对选定的诱导格给出了实际的模同构
\[
	\bigl(\mathcal OG\otimes_{\mathcal OH}L\bigr)/\pi\bigl(\mathcal OG\otimes_{\mathcal OH}L\bigr) \cong kG\otimes_{kH}(L/\pi L).
\]
而命题中关于 \(d_G,d_H\) 的公式位于 \(G_0\) 中：若先在特征零中分解诱导模，再为各直和项另外选格，其约化只保证具有相同的合成因子，未必与这个诱导格的约化同构。

\begin{example}\label{b5:exm:s3-induced-lattice}
	设 \(k\) 为特征 \(3\) 的域，\(G=S_3\)，\(H=\langle(12)\rangle\)，令 \(S_+,S_-\) 分别为平凡模与符号模。则置换模
	\[
		\operatorname{Ind}_H^G k\cong k^3
	\]
	不可分解且投射，合成因子为两个 \(S_+\) 与一个 \(S_-\)。令 \(U=\{(x,y,z)\in k^3:x+y+z=0\}\)，则 \(k^3\) 与 \(S_+\oplus U\) 的合成因子相同，却不同构。
\end{example}
\begin{proof}
	三个陪集上的置换作用就是三个坐标上的作用。限制到 \(C_3=\langle r\rangle\)、\(r=(123)\) 后，三个坐标构成一个正则轨道，故这个限制模是 \(kC_3\cong k[t]/(t^3)\)，其中 \(t=r-1\)。由\cref{b5:thm:cyclic-p-modules}它不可分解；若 \(k^3\) 有非平凡 \(kG\)-模直和分解，限制后也会分解，矛盾。 因为 \(|H|=2\) 在 \(k\) 中可逆，\(k\) 是投射 \(kH\)-模；诱导把自由模的直和因子送到自由模的直和因子，故 \(k^3\) 投射。

	在 \(U\) 中，\(D=k(1,1,1)\) 是平凡子模，\(U/D\) 是符号模：例如换位 \((12)\) 将 \((1,-1,0)\) 的类变为其负，而这个类非零。求和映射又给出 \(k^3/U\cong S_+\)，得到所列合成因子。\(r\) 在 \(U\) 上不是恒等，所以 \(0\to S_+\to U\to S_-\to0\) 不分裂；否则 \(r\) 会在两个直和项上都作用为恒等。\(S_+\oplus U\) 则有显示出的非平凡直和分解，因而不与 \(k^3\) 同构。
\end{proof}
在一个对 \(S_3,H\) 分裂的 \(3\)-模系统中，诱导格 \(\mathcal O[G/H]\) 的约化就是上述 \(k^3\)。特征零中由普通特征标或坐标求和可得
\[
	\operatorname{Ind}_H^G K\cong K\oplus U_K,\qquad U_K=\{(x,y,z)\in K^3:x+y+z=0\}.
\]
若为右侧另选直和格 \(\mathcal O\oplus\{(x,y,z)\in\mathcal O^3:x+y+z=0\}\)，其约化便是 \(S_+\oplus U\)。这两个格都位于同一个特征零表示中，分解映射将它们识别为同一个合成类，却没有将它们的约化模识别为同一个模。

\subsection{缺陷群与块论简介}
\label{b5:subsec:0804:02}

设 \(G\) 为有限群，\(p\) 为素数，\(k\) 为特征 \(p\) 的域，\(P\le G\) 为 \(p\)-子群。先把 \(P\) 的共轭作用当作置换作用来看。固定空间 \((kG)^P\) 以各个共轭轨道的轨道和为基，其中单点轨道恰好是 \(C_G(P)\) 中的元素。若一个轨道含 \(g\)，令 \(Q=C_P(g)\) 为它在 \(P\) 中的稳定子群，则其轨道和为
\[
	\sum_{x\in P/Q}xgx^{-1}.
\]
这里 \(g\in(kG)^Q\)。对任意子群 \(Q\le P\)，相对迹
\[
	\operatorname{Tr}_Q^P:(kG)^Q\longrightarrow(kG)^P,\qquad z\longmapsto\sum_{x\in P/Q}xzx^{-1}
\]
与代表元选择无关，因为 \(z\) 被 \(Q\) 固定；所得和被 \(P\) 固定。非单点轨道的稳定子群 \(Q\) 是真子群，所以该轨道和正是从较小子群来的相对迹。

\ProseParagraphBreak

反过来，若 \(Q<P\)，则 \(\operatorname{Tr}_Q^P(z)\) 在每个 \(h\in C_G(P)\) 上的系数为 \([P:Q]\) 乘以 \(z\) 在 \(h\) 上的系数；这个指数是正的 \(p\) 次幂，在 \(k\) 中为零。因此保留单点轨道、删去其余轨道的操作，其核恰为全部真子群相对迹的像之和。它保留的是固定空间中不能由这些较小子群的相对迹表示的部分：
\[
	(kG)^P\big/\sum_{Q<P} \operatorname{Tr}_Q^P\bigl((kG)^Q\bigr) \cong kC_G(P)
\]
作为 \(k\) 向量空间同构。下面还要验证，这个系数截断与乘法相容。

\begin{proposition}\label{b5:prop:brauer-truncation}
	设 \(G\) 为有限群，\(p\) 为素数，\(k\) 为特征 \(p\) 的域，\(P\le G\) 为 \(p\)-子群，\(C_G(P)\) 为其中心化子。令 \(P\) 以共轭作用于 \(kG\)，并记固定子代数为 \((kG)^P\)。映射
	\[
		\operatorname{Br}_P^G:(kG)^P\longrightarrow kC_G(P),\qquad \sum_{g\in G}a_gg\longmapsto\sum_{g\in C_G(P)}a_gg
	\]
	是保持单位的 \(k\) 代数同态，称为此处的\term[Brauer-yingshe]{Brauer 映射}[Brauer map]。
\end{proposition}
\begin{proof}
	加法、标量与单位的保持直接来自定义。取两个 \(P\) 共轭不变的元素 \(x=\sum a_gg\)、\(y=\sum b_hh\)，固定 \(z\in C_G(P)\)。乘积中 \(z\) 的系数为
	\[
		\sum_{gh=z}a_gb_h.
	\]
	\(P\) 同时共轭二元组 \((g,h)\)，保持条件 \(gh=z\)，且每个轨道内的系数乘积相同。非平凡轨道的大小是正的 \(p\) 次幂，在 \(k\) 中贡献为零。剩余的固定二元组恰好满足 \(g,h\in C_G(P)\)，其总贡献就是 \(\operatorname{Br}_P^G(x)\operatorname{Br}_P^G(y)\) 的 \(z\) 系数。逐个 \(z\) 比较即得乘法性。
\end{proof}
\begin{definition}\label{b5:def:block-defect-group}
	设 \(G\) 为有限群，\(p\) 为素数，\(k\) 为特征 \(p\) 的域，\(b\) 为 \(kG\) 的块幂等元，对 \(p\)-子群 \(D\le G\)，记 \(\operatorname{Br}_D^G:(kG)^D\to kC_G(D)\) 为保留中心化子上系数的 Brauer 映射。若 \(D\) 满足 \(\operatorname{Br}_D^G(b)\ne0\)，并且在满足这一条件的 \(p\)-子群中按包含关系极大，就称 \(D\) 为块 \(kGb\) 的一个\term[quexian-qun]{缺陷群}[defect group]。
\end{definition}
\symidx{\(\operatorname{Br}_D^G\)：群代数的 Brauer 映射}
由于 \(\operatorname{Br}_{1}^G(b)=b\ne0\)，且 \(G\) 有限，这样的极大子群存在。其进一步的共轭与结构性质属于块论；这里采用支集截断的描述以便计算。这个描述及其与一般定义的关系，可参阅\cite[Corollaries 12.5.4--12.5.5]{Webb2016FiniteGroupRepresentations}。
\ProseParagraphBreak
对\cref{b5:exm:s3-two-blocks}中的 \(S_3\) 与特征二，取 \(P=\langle s\rangle\)。中心化子 \(C_G(P)=P\)，故
\[
	\operatorname{Tr}_{1}^{P}(r)=r+srs=r+r^2.
\]
这是非单点共轭轨道的轨道和，Brauer 映射将它删去；单位元则是单点轨道，被保留下来。因此块幂等元 \(e=1+r+r^2\)、\(f=r+r^2\) 分别满足
\[
	\operatorname{Br}_P^G(1+r+r^2)=1,\qquad \operatorname{Br}_P^G(r+r^2)=0.
\]
三个非平凡 \(2\) 子群互相共轭，因此二维局部块的缺陷群为一个 Sylow \(2\) 子群，四维矩阵块的缺陷群为平凡群。后者称为缺陷零块。前者的简单模 \(S_1\) 不投射，其覆盖 \(kGe\) 有两个 \(S_1\) 合成因子，Cartan 条目为 \(2\)；后者的二维简单模 \(S_2\) 本身就投射，Cartan 条目为 \(1\)。这是本例中由块结构直接算出的对应，不代替一般缺陷零块理论的证明。在任意有限 \(p\) 群上，唯一块的幂等元为一，而对整个群的截断仍为一，所以其缺陷群就是整个群。

\ProseParagraphBreak

定义域中的共轭不变性不能省去。若把同样的支集截断定义在整个 \(kG\) 上，\(r,r^{-1}\) 都被送到零，但 \(rr^{-1}=1\) 被送到一，便不保持乘法。这两个群元素各自都不被 \(P\) 的共轭作用固定；必须先把它们组成轨道和，才能应用\cref{b5:prop:brauer-truncation}。

\subsection{特征零与正特征下的表示分类}
\label{b5:subsec:0804:03}

不可约 Brauer 特征标分类的是简单模，并由加法性确定任意模的合成因子。对于扩张信息，它有明确局限：循环 \(p\) 群上，\(M_d\) 与 \(k^{\oplus d}\) 的 Brauer 特征标都只在单位元处取整数 \(d\)，但前者不可分解，后者在 \(d>1\) 时可分解。分解矩阵同样保留约化后的合成因子，不指定某个整格约化的全部模结构。
\ProseParagraphBreak
有限群在正特征中甚至可以有无限多个不可分解表示。这里给一个不用分类定理的例子。设 \(k\) 是特征二的无限域，\(G=C_2\times C_2=\langle x,y\rangle\)，并在 \(k^2\) 上取
\[
	N=\begin{pmatrix}0&1\\0&0\end{pmatrix},\qquad
	\rho_\lambda(x)=I+N,\qquad \rho_\lambda(y)=I+\lambda N \quad(\lambda\in k).
\]
因为 \(N^2=0\)，两个矩阵交换且平方都为一，故给出表示。其自同态必须与 \(N\) 交换，恰为 \(aI+bN\)，组成局部环，所以每个表示不可分解。若两个参数的表示同构，同构矩阵 \(T\) 与 \(N\) 交换，再由 \(T(\lambda N)=\mu NT\) 得到 \((\lambda-\mu)NT=0\)。\(T\) 可逆且 \(N\ne0\)，故 \(\lambda=\mu\)。于是得到一个参数化的无限族。 记这些模为 \(V_\lambda\)，取基 \(e_1,e_2\) 使 \(Ne_2=e_1\)。由\cref{b5:prop:p-group-local}，\(A=kG\) 的根为 \(J=(x-1,y-1)\)，故对每个参数都有
\[
	JV_\lambda=NV_\lambda=ke_1,\qquad J^2V_\lambda=0, \qquad V_\lambda/JV_\lambda\cong k, \qquad \operatorname{soc}(V_\lambda)=ke_1\cong k.
\]
所以连全部根层也相同。根层各自只记录一个平凡模，而 \(y-1\) 在原模上是 \(x-1\) 的 \(\lambda\) 倍；这一比例记录两层之间的作用关系，却不出现在层商中。
\ProseParagraphBreak
稳定范畴把经过投射模分解的同态视为零。在循环 \(p\) 群上，这种商与扩张的联系也能直接算出。对有限维 \(k\) 代数 \(A\) 及其有限维左模 \(M,N\)，这里以 \(\operatorname{Ext}_A^1(M,N)\) 表示短正合列 \(0\to N\to E\to M\to0\) 的等价类：等价要求中项同构在两端诱导恒等，零类为分裂列。两个扩张先沿对角映射 \(M\to M\oplus M\) 拉回，再沿相加映射 \(N\oplus N\to N\) 推出，得到 Baer 和；沿 \(c1_N\) 推出给出标量 \(c\in k\) 的作用。下面的生成元描述将这些运算化为向量空间中的运算，不要求先掌握消解理论。

\begin{proposition}\label{b5:prop:cyclic-stable-hom-extensions}
	设 \(p\) 为素数，\(k\) 为特征 \(p\) 的域，\(r\ge1\) 为整数、\(m=p^r\)，令 \(A=k[t]/(t^m)\)、\(M_d=A/(t^d)\)，其中 \(1\le d\le m\)。对有限维左 \(A\)-模 \(M,N\)，记 \(\underline{\operatorname{Hom}}_A(M,N)\) 为 \(\operatorname{Hom}_A(M,N)\) 模掉所有可经过有限维投射 \(A\)-模分解的同态所得的商空间。对 \(1\le a,b\le m\)，有
	\[
		\underline{\operatorname{Hom}}_A(M_a,M_b) \cong\ker(t^a:M_b\to M_b)/t^{m-a}M_b.
	\]
	若 \(a<m\)，取投射覆盖 \(A\to M_a\) 的核为第一\term[hechongmo]{合冲模}[syzygy] \(\Omega M_a\)，则
	\[
		\begin{aligned}
			\Omega M_a&=t^aA\cong M_{m-a},\\
			\operatorname{Ext}_A^1(M_a,M_b) &\cong\underline{\operatorname{Hom}}_A(\Omega M_a,M_b)\\
			&\cong\ker(t^{m-a}:M_b\to M_b)/t^aM_b.
		\end{aligned}
	\]
	稳定同态空间与扩张空间的维数均为
	\[
		\min(a,b,m-a,m-b).
	\]
	\(a=m\) 时 \(M_a=A\) 投射，所有上述扩张分裂，维数仍由该式给出。
\end{proposition}
\symidx{\(\underline{\operatorname{Hom}}_A(M,N)\)：模掉经过投射模的同态所得的稳定同态空间}
\symidx{\(\Omega M\)：投射覆盖的核，即第一合冲模}
\symidx{\(\operatorname{Ext}_A^1(M,N)\)：固定两端的短正合模扩张类空间}
\begin{proof}
	同态 \(M_a\to M_b\) 由生成元一的像决定，该像恰好遍历 \(\ker(t^a:M_b\to M_b)\)。由\cref{b5:thm:split-projective-covers}，有限维投射 \(A\)-模都是有限个 \(A\) 的直和。\(M_a\to A\) 中一的像必须属于 \(t^{m-a}A\)，再经 \(A\to M_b\)，所得像恰好遍历 \(t^{m-a}M_b\)。这已是一个子空间，经过若干份 \(A\) 的映射之和不会再扩大它，得到第一式。核的维数为 \(\min(a,b)\)，分母的维数为 \(\max(0,a+b-m)\)，若 \(a+b\le m\)，则 \(m-a\ge b\)、\(m-b\ge a\)，两者之差为 \(\min(a,b)=\min(a,b,m-a,m-b)\)。若 \(a+b>m\)，可由对称性设 \(a\le b\)，差便为 \(a-(a+b-m)=m-b\)；此时 \(m-b<a\le b\)，且 \(m-b\le m-a\)，所以差仍为四者的最小值。

	对 \(a<m\)，映射 \(M_{m-a}\to t^aA\)、\(1\mapsto t^a\) 是同构。给定扩张，取商模生成元的一个提升 \(v\in E\)。因为商上 \(t^a=0\)，有 \(t^av=w\in M_b\)；又因 \(t^mE=0\)，必有 \(t^{m-a}w=0\)。反过来，给定这样的 \(w\)，在
	\[
		M_b\oplus\operatorname{span}_k\{v,tv,\ldots,t^{a-1}v\}
	\]
	上令 \(t\) 在第一项按原作用，在第二组向量上逐次后移，最后把 \(t^{a-1}v\) 送到 \(w\)。条件 \(t^{m-a}w=0\) 保证 \(t^m=0\)，所以这定义了一个具有指定核与商的 \(A\)-模扩张。 改变提升为 \(v+z\)、\(z\in M_b\)，会把参数改为 \(w+t^az\)；固定两端的扩张同构也只能这样改变提升。因此扩张类恰由 \(\ker(t^{m-a})/t^aM_b\) 参数化。Baer 和的拉回中可以取两个提升组成的向量，随后推出把两个核中的参数相加；标量推出则把参数乘以同一标量。故这是向量空间同构，零类恰为分裂扩张。

	将第一式中的 \(a\) 换成 \(m-a\)，便得到这个商空间正是 \(\underline{\operatorname{Hom}}_A(\Omega M_a,M_b)\)。也可从投射覆盖的短正合列理解分母：\(\Omega M_a\) 的生成元是 \(t^a\)，所以限制一个 \(A\to M_b\) 的同态时，该生成元的像恰好遍历 \(t^aM_b\)。扩张空间的维数为
	\[
		\min(m-a,b)-\max(0,b-a)=\min(a,b,m-a,m-b).
	\]
	若 \(a=m\)，投射性使每个到 \(M_a\) 的满射分裂，最后的断言也成立。
\end{proof}
所以这两个空间维数相同并非偶合，但一般并没有在这里把 \(\operatorname{Ext}_A^1(M_a,M_b)\) 与 \(\underline{\operatorname{Hom}}_A(M_a,M_b)\) 识别。与扩张自然相连的是第一变量换成合冲模以后的稳定同态空间。

\ProseParagraphBreak

这些计算解释了为何后续会研究投射覆盖、块、相对投射性和稳定范畴，而不只寻找另一张有限的特征标表。一般缺陷群如何控制块、何时只有有限多种不可分解模，以及不同块之间的等价，都是进一步的问题。本编接下来改从箭图研究有限维代数；那里的表示可以在任意特征上定义，所用的无环性、关系和代数闭性条件仍须逐项写明。
